\section{Thesis structure}\label{section:intro-thesis-structure}

The remainder of this thesis is structured as follows. 

Chapter~\ref{chapter:Related_Work} reviews existing research on scientific document extraction, biomedical information extraction and entity normalization, LLM-based structured extraction, natural language inference, provenance, cost-aware model cascading, and the evaluation of LLM-based extraction systems, then explains what the thesis adopts from prior work and where it diverges.

Chapter~\ref{chapter:Methodology} describes the system and the evaluation design. It presents the two-pipeline architecture, the corpus and how it was assembled, the document-extraction pipeline, the knowledge-extraction pipeline and its typed MAP-schema, the database backbone that connects the two pipelines, the agreement-based cascade, and the evaluation methodology that defines how the five research questions are measured. 

Chapter~\ref{chapter:Results} reports the empirical findings in the order of the research questions, from document extraction to held-out generalization.

Chapter~\ref{chapter:Discussion} interprets the results without introducing new measurements and discusses the main threats to validity. 

Chapter~\ref{chapter:Conclusion_and_Future_Work} summarizes the contributions and outlines future work, including human-annotated gold labels, generalization tests beyond the current corpus, and relation classification on real pipeline outputs, rather than on a synthetic dataset. 